\chapter{Continued Fractions and Pell's Equation}

\begin{center}\itshape
$\sqrt2=1+\cfrac1{2+\cfrac1{2+\cfrac1{2+\cdots}}}$, and the calendar's leap years come from the expansion of $0.24219$.\\
Continued fractions are Euclid's algorithm run on a real number, and they find the best rational approximations.
\end{center}

According to an epigram attributed to Archimedes, the Sun god's cattle grazed on the plains of Sicily, bulls and cows, white, black, yellow and dappled,
in proportions given by nine conditions, and the poem promises the solver perfect wisdom (we paraphrase). The problem reduces to a
\emph{Pell equation}, $x^2-4729494y^2=1$, together with additional divisibility conditions. Those extra conditions force a high power of a Pell unit: the smallest admissible herd has $206\,545$ digits, first computed in full in 1965. It is not simply the fundamental solution of the displayed Pell equation. Pell
equations are solved by continued fractions, which are Euclid's algorithm run on a real number instead of a fraction. This chapter develops them, and
shows that they give the best rational approximations of any real number, a fact used daily by gear designers, calendar makers and musicians.

\section{Finite continued fractions}

For real $a_0$ and $a_1,\dots,a_n>0$, define $[a_0]=a_0$ and $[a_0;a_1,\dots,a_n]=a_0+\dfrac1{[a_1;a_2,\dots,a_n]}$. Define the \emph{convergent numerators and denominators} by
\[
\begin{gathered}p_{-1}=1,\ p_0=a_0,\ p_k=a_kp_{k-1}+p_{k-2};\qq q_{-1}=0,\ q_0=1,\ q_k=a_kq_{k-1}+q_{k-2}.\end{gathered}
\]

\begin{proposition}\label{prop:cfconv}
\begin{enumerate}[label=(\alph*)]
\item For $n\ge1$ and every real $t>0$,
\[
[a_0;a_1,\dots,a_{n-1},t]=\frac{tp_{n-1}+p_{n-2}}{tq_{n-1}+q_{n-2}}.
\]
In particular $[a_0;\dots,a_n]=\frac{p_n}{q_n}$.
\item $p_nq_{n-1}-p_{n-1}q_n=(-1)^{n-1}$ for $n\ge0$. Hence, when the $a_i$ are integers, $\gcd(p_n,q_n)=1$, and
\[
\frac{p_n}{q_n}-\frac{p_{n-1}}{q_{n-1}}=\frac{(-1)^{n-1}}{q_nq_{n-1}}\qquad(n\ge1).
\]
\item If the $a_i$ ($i\ge1$) are positive integers, $q_n\ge F_{n+1}$ (Fibonacci), so $q_n\to\infty$.
\end{enumerate}
\end{proposition}
\begin{proof}
(a) Induction on $n$. For $n=1$, $[a_0;t]=a_0+\frac1t=\frac{ta_0+1}{t}=\frac{tp_0+p_{-1}}{tq_0+q_{-1}}$. For $n\ge2$, $[a_0;\dots,a_{n-1},t]=[a_0;\dots,a_{n-2},a_{n-1}+\frac1t]$, which by the induction hypothesis is
\[
\begin{aligned}\frac{(a_{n-1}+\frac1t)p_{n-2}+p_{n-3}}{(a_{n-1}+\frac1t)q_{n-2}+q_{n-3}}&=\frac{p_{n-1}+\frac1tp_{n-2}}{q_{n-1}+\frac1tq_{n-2}}\eqbr=\frac{tp_{n-1}+p_{n-2}}{tq_{n-1}+q_{n-2}}.\end{aligned}
\]
With $t=a_n$ the right side is $\frac{p_n}{q_n}$.

(b) For $n=0$: $p_0q_{-1}-p_{-1}q_0=-1$. Then $p_nq_{n-1}-p_{n-1}q_n=(a_np_{n-1}+p_{n-2})q_{n-1}-p_{n-1}(a_nq_{n-1}+q_{n-2})=-(p_{n-1}q_{n-2}-p_{n-2}q_{n-1})$. A common divisor of $p_n,q_n$ divides $\pm1$. Divide by $q_nq_{n-1}$.

(c) $q_0=1=F_1$, $q_1=a_1\ge1=F_2$, and $q_n\ge q_{n-1}+q_{n-2}$.
\end{proof}

For a rational $\frac ab$ with $b\ge1$, Euclid's algorithm $a=a_0b+r_1$, $b=a_1r_1+r_2,\dots$ gives $\frac ab=a_0+\frac{r_1}b=a_0+\cfrac1{b/r_1}=\cdots=[a_0;a_1,\dots,a_n]$: the partial quotients of a rational number are the quotients of Euclid's algorithm.

\subsection*{A fraction dismantled, one floor at a time}
Take $43/19$. First remove its integer part, $2$; then turn the remainder $5/19$ upside down. Repeat:
\[
\frac{43}{19}=2+\frac5{19},\qquad
\frac{19}5=3+\frac45,\qquad
\frac54=1+\frac14.
\]
Thus $43/19=[2;3,1,4]$. The recurrence is a bookkeeping device: starting with $(p_{-1},q_{-1})=(1,0)$ and $(p_0,q_0)=(2,1)$ gives
\[
\frac21,\qquad\frac73,\qquad\frac94,\qquad\frac{43}{19}.
\]
At the last step $4\cdot9+7=43$ and $4\cdot4+3=19$. We retain only the previous two columns, however long the expansion becomes.

The last partial quotient of a non-integral rational may be taken at least $2$. This convention removes the harmless ambiguity
\[
[2;3,1,4]=[2;3,1,3,1],
\]
because $3+1/1=4$. Infinite expansions have no final digit to split this way. Also, $a_0$ may be negative: the algorithm uses the \emph{floor}, so $-7/3=[-3;1,2]$, not an integer part obtained by rounding toward zero.

\section{The expansion of a real number}

Let $x\in\R$. Put $x_0=x$ and set $a_k=\lfloor x_k\rfloor$ at each step. If $x_k$ is an integer, stop; otherwise put $x_{k+1}=\frac1{x_k-a_k}>1$. Then $a_k\ge1$ for $k\ge1$ and $x=[a_0;a_1,\dots,a_{k-1},x_k]$ (induction, since $x_k=a_k+\frac1{x_{k+1}}$). The $x_k$ are the \emph{complete quotients}. For rational $x$ the process stops (it is Euclid's algorithm); for irrational $x$ it never stops, since every $x_k$ is irrational.

\begin{think}
Take $\sqrt2\approx1.41421$ and its first fractions $1,\frac32,\frac75,\frac{17}{12},\frac{41}{29}$. Which are below $\sqrt2$, and which above? Multiply each error $\frac pq-\sqrt2$ by $q^2$. What do you notice, and what would you guess for the next one?
\end{think}

\begin{theorem}\label{thm:cfapprox}
Let $x$ be irrational with partial quotients $a_k$ and convergents $\frac{p_k}{q_k}$. Then for every $n\ge0$,
\[
\begin{gathered}x-\frac{p_n}{q_n}=\frac{(-1)^n}{q_n\left(x_{n+1}q_n+q_{n-1}\right)},\qq\text{so}\qq\frac1{q_n(q_n+q_{n+1})}<\left|x-\frac{p_n}{q_n}\right|<\frac1{q_nq_{n+1}}\le\frac1{q_n^2}.\end{gathered}
\]
In particular $\frac{p_n}{q_n}\to x$, the convergents lie alternately below and above $x$, and we write $x=[a_0;a_1,a_2,\dots]$.
\end{theorem}
\begin{proof}
Let $Q=x_{n+1}q_n+q_{n-1}$. By Proposition~\ref{prop:cfconv}(a) with $t=x_{n+1}$, $x=\frac{x_{n+1}p_n+p_{n-1}}Q$, so
\[
\begin{aligned}x-\frac{p_n}{q_n}&=\frac{q_n(x_{n+1}p_n+p_{n-1})-p_nQ}{q_nQ}\\&=\frac{p_{n-1}q_n-p_nq_{n-1}}{q_nQ}=\frac{(-1)^n}{q_nQ}.\end{aligned}
\]
Since $a_{n+1}<x_{n+1}<a_{n+1}+1$, the number $Q$ lies strictly between $q_{n+1}$ and $q_{n+1}+q_n$. Finally $q_n\to\infty$.
\end{proof}

Conversely, every sequence $a_0\in\Z$, $a_1,a_2,\dots\ge1$ is the expansion of exactly one irrational number (Exercise~\ref{ex:cfunique}).

\section{Best approximations}

\begin{think}
Among all fractions with denominator at most $10$, which is closest to $\pi$? And with denominator at most $100$? (A short program helps for the second.) Compare your answers with the convergents $3,\frac{22}7,\frac{333}{106},\frac{355}{113}$ of $\pi$. One of your answers is not a convergent: can you see why this does not contradict the lemma below, which measures $|qx-p|$ rather than $\left|x-\frac pq\right|$?
\end{think}

\begin{lemma}[best approximation]\label{lem:best}
Let $x$ be irrational and $n\ge0$. If $p,q$ are integers with $1\le q<q_{n+1}$, then $|qx-p|\ge|q_nx-p_n|$.
\end{lemma}
\begin{proof}
The matrix $\begin{pmatrix}p_n&p_{n+1}\\q_n&q_{n+1}\end{pmatrix}$ has determinant $\pm1$, so there are integers $u,v$ with $p=up_n+vp_{n+1}$ and $q=uq_n+vq_{n+1}$. If $u=0$, then $q=vq_{n+1}$ with $v\ne0$, so $q\ge q_{n+1}$ (as $q>0$): excluded. If $v=0$, $|qx-p|=|u|\,|q_nx-p_n|\ge|q_nx-p_n|$. Otherwise $u$ and $v$ have opposite signs (if both were $>0$, $q>q_{n+1}$; if both $<0$, $q<0$). By Theorem~\ref{thm:cfapprox}, $q_nx-p_n$ and $q_{n+1}x-p_{n+1}$ also have opposite signs. So $u(q_nx-p_n)$ and $v(q_{n+1}x-p_{n+1})$ have the same sign, and
$|qx-p|=|u(q_nx-p_n)|+|v(q_{n+1}x-p_{n+1})|\ge|q_nx-p_n|$.
\end{proof}

\begin{theorem}[Legendre]\label{thm:legendre}
Let $x$ be irrational, $p\in\Z$, $q\ge1$. If $\left|x-\frac pq\right|<\frac1{2q^2}$, then $\frac pq$ is a convergent of $x$.
\end{theorem}
\begin{proof}
Let $n$ be such that $q_n\le q<q_{n+1}$ (choose the largest index with $q_n\le q$; $q_0=q_1=1$ is possible, but thereafter the denominators strictly increase to $\infty$). By the lemma, $|q_nx-p_n|\le|qx-p|<\frac1{2q}$. Suppose $\frac pq\ne\frac{p_n}{q_n}$. Then $|pq_n-p_nq|\ge1$, so
\[
\begin{aligned}\frac1{qq_n}\le\left|\frac pq-\frac{p_n}{q_n}\right|&\le\left|x-\frac pq\right|+\left|x-\frac{p_n}{q_n}\right|\eqbr<\frac1{2q^2}+\frac1{2qq_n}.\end{aligned}
\]
Multiplying by $2qq_n$: $2<\frac{q_n}q+1$, so $q<q_n$, a contradiction.
\end{proof}

\subsection*{Two different meanings of ``best''}
The quantities
\[
\left|x-\frac pq\right|\quad\text{and}\quad|qx-p|=q\left|x-\frac pq\right|
\]
ask different questions. The first measures the error in the number itself; the second measures the miss when $q$ copies of $x$ are compared with an integer. The lemma proves a record for the second quantity. Dividing by different denominators can change the winner for the first.

For $\pi$, among denominators at most $100$, $311/99$ is closer than $22/7$ in the first sense:
\[
\begin{aligned}
\left|\pi-\frac{311}{99}\right|&\approx0.0001785,
&\left|\pi-\frac{22}7\right|&\approx0.0012645.
\end{aligned}
\]
But $|99\pi-311|\approx0.01767$ is \emph{larger} than $|7\pi-22|\approx0.00885$, exactly as the lemma requires. The number $311/99$ is an intermediate fraction, obtained from $(14\cdot22+3)/(14\cdot7+1)$. Its existence is a reason to read a theorem's metric as carefully as its conclusion.

\begin{figure}[ht]
\centering
\fig{ch08-lattice-errors}
\caption{Approximation as a lattice problem. A point $(q,p)$ represents $p/q$; its vertical distance from $p=\sqrt2\,q$ is $|q\sqrt2-p|$, while its slope error is $|\sqrt2-p/q|$. The labelled convergents alternate below and above the line. Distances, rather than colours, tell the story.}
\label{fig:cflattice}
\end{figure}

\begin{reminder}[A very good approximation is not an equality test]
Legendre's hypothesis is a sufficient condition: error less than $1/(2q^2)$ guarantees a convergent. It is not necessary for every convergent. Nor does seeing a long decimal agreement establish irrationality or a special identity. Continued fractions turn guesses into certified bounds; a calculator supplies the guesses.
\end{reminder}

So the convergents include all approximations satisfying Legendre's bound, and every irrational $x$ has infinitely many $\frac pq$ with $\left|x-\frac pq\right|<\frac1{q^2}$ (\emph{Dirichlet's theorem}). For rational $x=\frac ab$ this fails: $\left|\frac ab-\frac pq\right|\ge\frac1{bq}$ when $\frac pq\ne\frac ab$. This gives a criterion of irrationality that we use with $e$ numerically.

Figure~\ref{fig:quality} compares the quantity $q_n^2\left|x-\frac{p_n}{q_n}\right|$, which Theorem~\ref{thm:cfapprox} bounds by $1$, for four famous numbers. A large partial quotient $a_{n+1}$ makes it small, so the convergent $\frac{p_n}{q_n}$ is exceptionally good.

\begin{figure}[ht]
\centering
\fig{ch08-quality}
\caption{How good are the convergents? For $\varphi$ and $\sqrt2$, whose partial quotients are constant, $q_n^2\left|x-\frac{p_n}{q_n}\right|$ settles at $\frac1{\sqrt5}\approx0.447$ and $\frac1{2\sqrt2}\approx0.354$. For $\pi$, the dip at $n=3$ is $\frac{355}{113}$, followed by the huge partial quotient $292$: $113^2\left|\pi-\frac{355}{113}\right|\approx0.0034$. For $e$, every third value dips, following the pattern $1,2k,1$ of its partial quotients.}
\label{fig:quality}
\end{figure}

\begin{remark}
Hurwitz (1891) improved $\frac1{q^2}$ to $\frac1{\sqrt5\,q^2}$, and $\sqrt5$ is optimal because of the golden ratio (Exercise~\ref{ex:golden}). The proof is short: of any two consecutive convergents one satisfies $|x-\frac pq|<\frac1{2q^2}$, and of any three, one satisfies $|x-\frac pq|<\frac1{\sqrt5\,q^2}$ (see Hardy and Wright, Chapter~11, in the sources).
\end{remark}

\section{Periodic expansions}

A \emph{quadratic irrational} is an irrational root of a polynomial $AX^2+BX+C$ with $A,B,C\in\Z$, $A\ne0$.

\begin{theorem}[Euler, Lagrange]\label{thm:periodic}
The continued fraction of an irrational $x$ is eventually periodic if and only if $x$ is a quadratic irrational.
\end{theorem}
\begin{proof}
\emph{Periodic implies quadratic} (Euler). Suppose $a_{k+\ell}=a_k$ for all $k\ge m$. Then $x_{m+\ell}$ and $x_m$ have the same expansion, so $x_{m+\ell}=x_m=:y$ (Exercise~\ref{ex:cfunique}), and $y=[a_m;\dots,a_{m+\ell-1},y]=\frac{yP+P'}{yQ+Q'}$ with integers $P,P',Q,Q'$ by Proposition~\ref{prop:cfconv}(a). So $Qy^2+(Q'-P)y-P'=0$ with $Q\ge1$: $y$ is a quadratic irrational. The set $\Q(y)=\{r+sy:r,s\in\Q\}$ is a field: it is stable under products (reduce $y^2$ with the equation) and $\frac1{r+sy}$ is in it (multiply by the conjugate $r+sy'$, where $y'$ is the other root, and note $(r+sy)(r+sy')\in\Q^*$). If $m=0$, then $x=y$ and we are done. If $m\ge1$, $x=\frac{yp_{m-1}+p_{m-2}}{yq_{m-1}+q_{m-2}}$, so $x\in\Q(y)$, a $\Q$-vector space of dimension $2$; so $1,x,x^2$ are linearly dependent over $\Q$, and clearing denominators $x$ is a root of a non-zero integer polynomial of degree $\le2$, of degree exactly $2$ since $x$ is irrational.

\emph{Quadratic implies periodic} (Lagrange). Let $f=AX^2+BX+C$ vanish at $x$, with discriminant $D=B^2-4AC>0$, not a square ($x$ is irrational). For $n\ge2$, substitute $x=\frac{x_np_{n-1}+p_{n-2}}{x_nq_{n-1}+q_{n-2}}$ in $f(x)=0$ and multiply by $(x_nq_{n-1}+q_{n-2})^2$: $A_nx_n^2+B_nx_n+C_n=0$ with
\[
\begin{gathered}A_n=q_{n-1}^2f\!\left(\tfrac{p_{n-1}}{q_{n-1}}\right),\qq C_n=q_{n-2}^2f\!\left(\tfrac{p_{n-2}}{q_{n-2}}\right)=A_{n-1},\end{gathered}
\]
integers, and $B_n^2-4A_nC_n=D$. Indeed, write $f=A(X-x)(X-x')$, so $D=A^2(x-x')^2$, and $X=\frac{aY+b}{cY+e}$ with $ae-bc=\pm1$. Then $(cY+e)^2f\!\left(\frac{aY+b}{cY+e}\right)=A\bigl((a-cx)Y-(ex-b)\bigr)\bigl((a-cx')Y-(ex'-b)\bigr)=A'(Y-y)(Y-y')$ with $A'=A(a-cx)(a-cx')$, $y=\frac{ex-b}{a-cx}$, $y'=\frac{ex'-b}{a-cx'}$, and a direct computation gives $y-y'=\frac{(ae-bc)(x-x')}{(a-cx)(a-cx')}$, so $A'^2(y-y')^2=A^2(x-x')^2$. Now $A_n\ne0$ since $f$ has no rational root, and by Theorem~\ref{thm:cfapprox}, $p_{n-1}=xq_{n-1}+\frac{\delta}{q_{n-1}}$ with $|\delta|<1$, so
\[
\begin{gathered}\begin{aligned}A_n&=q_{n-1}^2f(x)+(2Ax+B)\delta+\frac{A\delta^2}{q_{n-1}^2}\eqbr=(2Ax+B)\delta+\frac{A\delta^2}{q_{n-1}^2},\end{aligned}\qq|A_n|<|2Ax+B|+|A|.\end{gathered}
\]
So $A_n$ and $C_n=A_{n-1}$ are bounded, and so is $B_n^2=D+4A_nC_n$. There are finitely many triples $(A_n,B_n,C_n)$, each with at most two roots, so the complete quotients $x_n$ take finitely many values: $x_m=x_{m+\ell}$ for some $m$ and $\ell\ge1$. Since each $x_{k+1}$ is computed from $x_k$ alone, $x_{k+\ell}=x_k$ for all $k\ge m$, and the partial quotients are periodic.
\end{proof}

\begin{example}
$\sqrt2=1+(\sqrt2-1)$ and $\frac1{\sqrt2-1}=\sqrt2+1=2+(\sqrt2-1)$, so $x_k=\sqrt2+1$ for all $k\ge1$ and $\sqrt2=[1;2,2,2,\dots]=[1;\overline2]$. Its convergents $1,\frac32,\frac75,\frac{17}{12},\frac{41}{29},\dots$ satisfy $p^2-2q^2=\pm1$.
\end{example}

\section{Pell's equation}

\begin{think}
Find by hand the least solution with $x,y\ge1$ of $x^2-2y^2=1$, of $x^2-3y^2=1$ and of $x^2-7y^2=1$. Now try $x^2-61y^2=1$: why does trying $y=1,2,3,\dots$ get you nowhere? The earlier thinking prompt about $\sqrt2$ suggests where to look.
\end{think}

Let $d\ge2$ be an integer that is not a square. \emph{Pell's equation} is $x^2-dy^2=1$ in integers. Write $N(x+y\sqrt d)=x^2-dy^2=(x+y\sqrt d)(x-y\sqrt d)$; it is multiplicative, because $(x+y\sqrt d)(x'+y'\sqrt d)=(xx'+dyy')+(xy'+x'y)\sqrt d$ and the same with $-\sqrt d$. And $x+y\sqrt d$ determines $(x,y)$, as $\sqrt d$ is irrational.

\begin{theorem}[Lagrange, 1768]\label{thm:pell}
Pell's equation has a solution with $y\ne0$. If $\varepsilon=x_1+y_1\sqrt d$ is the least solution with $x_1,y_1\ge1$ (the \emph{fundamental solution}), the solutions with $x,y\ge1$ are exactly the $x_k+y_k\sqrt d=\varepsilon^k$, $k\ge1$, and all solutions are $\pm\varepsilon^k$, $k\in\Z$.
\end{theorem}
\begin{proof}
\emph{Existence.} By Theorem~\ref{thm:cfapprox} there are infinitely many pairs $(p,q)$ with $q\ge1$ and $|p-q\sqrt d|<\frac1q$; for them $p>0$ and
\[
\begin{aligned}|p^2-dq^2|&=|p-q\sqrt d|\,|p+q\sqrt d|\eqbr<\frac1q\left(2q\sqrt d+\frac1q\right)\le2\sqrt d+1.\end{aligned}
\]
So some integer $k$ with $0<|k|\le2\sqrt d+1$ ($k\ne0$ as $\sqrt d\notin\Q$) is $p^2-dq^2$ for infinitely many such pairs. Among them, by pigeonhole on residues modulo $|k|$, there are two distinct pairs $(p,q)\ne(p',q')$ with $p\equiv p'$ and $q\equiv q'\pmod{|k|}$. Let $\alpha=p-q\sqrt d$ and $\alpha'=p'-q'\sqrt d$. Then
\[
\begin{aligned}\frac\alpha{\alpha'}&=\frac{(p-q\sqrt d)(p'+q'\sqrt d)}{k}\eqbr=\frac{(pp'-dqq')+(pq'-p'q)\sqrt d}k,\end{aligned}
\]
and $pp'-dqq'\equiv p^2-dq^2=k\equiv0$, $pq'-p'q\equiv pq-pq=0\pmod{|k|}$. So $\frac\alpha{\alpha'}=X+Y\sqrt d$ with $X,Y\in\Z$, and $X^2-dY^2=\frac{N(\alpha)}{N(\alpha')}=1$. If $Y=0$ then $X=\pm1$ and $\alpha=\pm\alpha'$, i.e.\ $(p,q)=\pm(p',q')$, impossible as the pairs are distinct with $p,p'>0$.

\emph{Structure.} If $u=x+y\sqrt d$ has norm $1$, so do $-u$, $u^{-1}=x-y\sqrt d$ and $-u^{-1}$; one of these four is $\ge1$. If $u>1$ has norm $1$, then $u^{-1}=x-y\sqrt d\in\left]0,1\right[$, so $x=\frac{u+u^{-1}}2>0$ and $y\sqrt d=\frac{u-u^{-1}}2>0$. So the solutions $>1$ are exactly those with $x,y\ge1$; there is one (existence). Choose a solution with the least positive integer $y$; since $x=\sqrt{1+dy^2}$ and $x+y\sqrt d$ both increase with $y$, this gives the least solution $\varepsilon>1$. For a solution $u>1$, let $k\ge0$ with $\varepsilon^k\le u<\varepsilon^{k+1}$. Then $u\varepsilon^{-k}$ has norm $1$ and lies in $[1,\varepsilon[$, so it equals $1$ by minimality: $u=\varepsilon^k$.
\end{proof}

\begin{example}
For $d=2$, $\varepsilon=3+2\sqrt2$ (try $y=1$: $x^2=3$, no; $y=2$: $x=3$), and the solutions are $(3,2),(17,12),(99,70),(577,408),\dots$ For $d=61$, the fundamental solution is $(1766319049,226153980)$, found by Bh\=askara~II in 1150; the continued fraction of $\sqrt{61}$, of period $11$, finds it at once.
\end{example}

\subsection*{How to find the fundamental solution without guessing}
Existence and structure leave an algorithmic question: where should we look for the first unit? If $x^2-dy^2=1$ with $x,y\ge1$, then
\[
0<\frac xy-\sqrt d
=\frac{1}{y^2(x/y+\sqrt d)}
<\frac1{2y^2},
\]
since $d\ge2$. Legendre's theorem therefore says that $x/y$ is a convergent of $\sqrt d$. Compute the convergents in order and test $p_n^2-dq_n^2=1$ using \emph{integer} arithmetic. Existence guarantees that the search ends; after the possible initial repeated denominator, the denominators grow strictly. Among positive Pell solutions $x=\sqrt{1+dy^2}$ grows with $y$, so the first success is the fundamental solution.

For $d=7$, the complete quotients can be simplified by rationalising:
\[
\begin{aligned}
\sqrt7&=2+\frac1{(\sqrt7+2)/3},\\
\frac{\sqrt7+2}{3}&=1+\frac1{(\sqrt7+1)/2},\\
\frac{\sqrt7+1}{2}&=1+\frac1{(\sqrt7+1)/3},\\
\frac{\sqrt7+1}{3}&=1+\frac1{\sqrt7+2}.
\end{aligned}
\]
The next integer part is $4$ and the cycle restarts, giving $\sqrt7=[2;\overline{1,1,1,4}]$. The convergents begin $2/1,3/1,5/2,8/3$; their norms are $-3,2,-3,1$. Thus $8+3\sqrt7$ is fundamental, and squaring gives the next solution $127+48\sqrt7$. A search through all $y$ would eventually find this too, but the convergents have already done the sifting.

\begin{think}
If $x^2-dy^2=-1$, why does squaring $x+y\sqrt d$ give a solution of norm $1$? Can the negative equation always be solved? For $d=3$, reducing modulo $3$ would require $x^2\equiv-1$, so the answer is no. The signs $+1$ and $-1$ are close on the page and quite different arithmetically.
\end{think}

\section{Exercises}

\begin{exercise}[$\star$]\label{ex:cf415}
Expand $\frac{415}{93}$ as a continued fraction, list its convergents, and use Proposition~\ref{prop:cfconv}(b) to solve $415x-93y=1$.
\end{exercise}

\begin{exercise}[$\star$]\label{ex:cfunique}
\begin{enumerate}[label=(\alph*)]
\item Let $a_0\in\Z$ and $a_1,a_2,\dots\ge1$ be integers, with convergents $\frac{p_n}{q_n}$. Show that $\frac{p_n}{q_n}$ converges to an irrational $x$ whose expansion is $[a_0;a_1,a_2,\dots]$.
\item Show that $[1;1,1,1,\dots]=\varphi=\frac{1+\sqrt5}2$, with convergents $\frac{F_{n+2}}{F_{n+1}}$, and that $\sqrt3=[1;\overline{1,2}]$.
\end{enumerate}
\end{exercise}

\begin{exercise}[$\star\star$]\label{ex:calendar}
The tropical year lasts about $365.24219$ days. Expand $0.24219$ as a continued fraction and list its first convergents. Interpret the Julian calendar (one leap year every $4$ years), the idealised $33$-year Persian cycle ($8$ leap years every $33$ years), and compare with the Gregorian rule ($97$ leap years every $400$).
\end{exercise}

\begin{exercise}[$\star\star$]\label{ex:n21}
Let $n\ge1$. Show that $\sqrt{n^2+1}=[n;\overline{2n}]$, and that $(2n^2+1,2n)$ is a solution of $x^2-(n^2+1)y^2=1$, obtained by squaring a solution of $x^2-(n^2+1)y^2=-1$.
\end{exercise}

\begin{exercise}[$\star\star$]\label{ex:sqtri}
Find all triangular numbers $\frac{m(m+1)}2$ that are perfect squares. \emph{Hint: reduce to $x^2-2y^2=1$.}
\end{exercise}

\begin{exercise}[$\star\star$]\label{ex:golden}
\emph{The golden ratio is the worst approximable number.} Show that for $x=\varphi$, $q_n^2\left|\varphi-\frac{p_n}{q_n}\right|\to\frac1{\sqrt5}$. Deduce that for every $c>\sqrt5$, there are only finitely many $\frac pq$ with $\left|\varphi-\frac pq\right|<\frac1{cq^2}$.
\end{exercise}

\begin{exercise}[$\star\star\star$]\label{ex:liouville}
\emph{Liouville's theorem (1844).} Let $x$ be a real root of an irreducible polynomial $P\in\Z[X]$ of degree $d\ge2$.
\begin{enumerate}[label=(\alph*)]
\item Show that there is $c>0$ with $\left|x-\frac pq\right|\ge\frac c{q^d}$ for all $p\in\Z$, $q\ge1$. \emph{Hint: $q^dP(p/q)$ is a non-zero integer; use the mean value theorem.}
\item Deduce that $L=\sum_{k\ge1}10^{-k!}$ is \emph{transcendental}: it is a root of no non-zero polynomial with integer coefficients.
\end{enumerate}
\end{exercise}

\begin{beauty}
The continued fraction of a number is a fingerprint: finite for rationals, periodic for quadratic irrationals, and for $e=[2;1,2,1,1,4,1,1,6,\dots]$ a strikingly regular pattern. For $\pi=[3;7,15,1,292,\dots]$ nobody sees any pattern, and the large $292$ explains why $\frac{355}{113}$ is so good. Pell's equation is our first glimpse of the \emph{units} of a ring of algebraic integers, the subject of Chapter~11 and of Dirichlet's unit theorem; Liouville's theorem gave the first explicit transcendental number, in 1844, before Hermite (1873) and Lindemann (1882) proved that $e$ and $\pi$ are transcendental.
\end{beauty}

\section{Hints}

\begin{hint}{ex:cf415}
Run Euclid's algorithm on $(415,93)$ and build the table of $p_k,q_k$. The last two convergents give $415q_{n-1}-93p_{n-1}=\pm1$.
\end{hint}
\begin{hint}{ex:cfunique}
(a) Even convergents increase and odd ones decrease, with gaps $\frac1{q_nq_{n+1}}$. Then show by induction that the complete quotients of the limit are $[a_k;a_{k+1},\dots]$, so the integer parts are the $a_k$. (b) $\varphi=1+\frac1\varphi$.
\end{hint}
\begin{hint}{ex:calendar}
The expansion is $0.24219=[0;4,7,1,3,\dots]$. Each convergent $\frac km$ reads as ``$k$ leap years every $m$ years''. Compare the yearly errors.
\end{hint}
\begin{hint}{ex:n21}
Rationalise $\frac1{\sqrt{n^2+1}-n}$. For Pell, $n^2-(n^2+1)\cdot1^2=-1$; square $n+\sqrt{n^2+1}$. \emph{If stuck:} $(a+b\sqrt d)^2=(a^2+db^2)+2ab\sqrt d$; identify the two coefficients before taking a norm.
\end{hint}
\begin{hint}{ex:sqtri}
Multiply $\frac{m(m+1)}2=k^2$ by $8$ and complete the square: $(2m+1)^2-2(2k)^2=1$. Then use Theorem~\ref{thm:pell} with $\varepsilon=3+2\sqrt2$, and check which solutions have $x$ odd and $y$ even.
\end{hint}
\begin{hint}{ex:golden}
Theorem~\ref{thm:cfapprox} gives the exact error $\frac1{q_n(x_{n+1}q_n+q_{n-1})}$. For $\varphi$, every $x_{n+1}=\varphi$ and $\frac{q_{n-1}}{q_n}\to\frac1\varphi$. For the second part, combine with Legendre's theorem.
\end{hint}
\begin{hint}{ex:liouville}
(a) Near $x$, $|P(x)-P(\frac pq)|\le M|x-\frac pq|$ with $M=\sup|P'|$ on $[x-1,x+1]$, and $|P(\frac pq)|\ge\frac1{q^d}$. (b) Truncate the series at $k=n$: the error is about $10^{-(n+1)!}$, far smaller than $\frac c{q^d}$ with $q=10^{n!}$. \emph{A second hint:} first prove that $L$ is irrational using the lower bound $1/(bq)$ for two distinct rationals with denominators $b,q$. Otherwise a hypothetical degree-$1$ polynomial would escape part (a).
\end{hint}

\section{Solutions}

\begin{solution}{ex:cf415}
\idea The partial quotients of a rational are the quotients of Euclid's algorithm, and the identity $p_nq_{n-1}-p_{n-1}q_n=\pm1$ applied to the last convergent $\frac{p_n}{q_n}=\frac{415}{93}$ is a B\'ezout relation, read off from the previous convergent.
\writeup Euclid: $415=4\cdot93+43$, $93=2\cdot43+7$, $43=6\cdot7+1$, $7=7\cdot1$. So $\frac{415}{93}=[4;2,6,7]$, with convergents $\frac41,\frac92,\frac{58}{13},\frac{415}{93}$ (from $p_k=a_kp_{k-1}+p_{k-2}$: $9=2\cdot4+1$, $58=6\cdot9+4$, $415=7\cdot58+9$, and $2,13=6\cdot2+1,93=7\cdot13+2$). By Proposition~\ref{prop:cfconv}(b) with $n=3$, $415\cdot13-58\cdot93=(-1)^2=1$. So $(x,y)=(13,58)$ is a solution, and all solutions are $(13+93k,58+415k)$, $k\in\Z$ (as in Exercise~\ref{ex:bez}).
\end{solution}

\begin{solution}{ex:cfunique}
\idea (a) The convergents form two adjacent sequences (even ones increase, odd ones decrease, gap $\frac1{q_nq_{n+1}}\to0$), hence converge. To show that the expansion of the limit is the given sequence, show that the expansion algorithm applied to $x$ recovers $a_0$, then shift. (b) For a periodic expansion, the tail equals the whole number, which gives an equation.
\writeup (a) By Proposition~\ref{prop:cfconv}(b), $\frac{p_{n+2}}{q_{n+2}}-\frac{p_n}{q_n}=\frac{(-1)^{n+1}}{q_{n+2}q_{n+1}}+\frac{(-1)^n}{q_{n+1}q_n}=(-1)^n\frac{q_{n+2}-q_n}{q_nq_{n+1}q_{n+2}}$, positive for even $n$ and negative for odd $n$. So the even convergents increase, the odd ones decrease, and their gap tends to $0$ ($q_n\to\infty$): they are adjacent sequences with a common limit $x$, and $\frac{p_{2m}}{q_{2m}}<x<\frac{p_{2m+1}}{q_{2m+1}}$. The same argument applies to every tail: let $y=\lim_n[a_1;a_2,\dots,a_n]$ and $z=\lim_n[a_2;a_3,\dots,a_n]$. Passing to the limit in $[a_0;a_1,\dots,a_n]=a_0+\frac1{[a_1;\dots,a_n]}$ gives $x=a_0+\frac1y$, and likewise $y=a_1+\frac1z$. The strict inequalities above give $z>[a_2]\ge1$ (the limit is strictly above the even convergents, the first of which is $a_2$), so $y>a_1\ge1$ and $0<\frac1y<1$. Hence $\lfloor x\rfloor=a_0$ and the first complete quotient of $x$ is $x_1=y$; by induction the whole expansion of $x$ is $[a_0;a_1,\dots]$. It is infinite, so $x$ is irrational. Two irrationals with the same expansion are the limit of the same convergents, hence equal.

(b) If $x=[1;1,1,\dots]$, the tail is $x$ again, so $x=1+\frac1x$, $x^2-x-1=0$, $x>1$: $x=\varphi$. The recurrences give $p_n=F_{n+2}$ and $q_n=F_{n+1}$. For $\sqrt3$: $\lfloor\sqrt3\rfloor=1$, $x_1=\frac1{\sqrt3-1}=\frac{\sqrt3+1}2$, so $a_1=1$; $x_2=\frac1{\frac{\sqrt3+1}2-1}=\frac2{\sqrt3-1}=\sqrt3+1$, so $a_2=2$ and $x_3=\frac1{\sqrt3-1}=x_1$. So $\sqrt3=[1;\overline{1,2}]$.
\end{solution}

\begin{solution}{ex:calendar}
\idea A calendar rule ``$k$ leap years every $m$ years'' approximates $0.24219$ by $\frac km$, and the convergents are the best such approximations for a given size of $m$ (Lemma~\ref{lem:best}). So expand, and compare the historical rules with the convergents.
\writeup Running the algorithm,
\[\begin{gathered}0.24219=[0;4,7,1,3,24,\dots],\qq\text{convergents }\frac14,\ \frac7{29},\ \frac8{33},\ \frac{31}{128},\ \frac{752}{3105},\ \dots\end{gathered}\]
The Julian calendar uses $\frac14=0.25$: too many leap years, by about $0.0078$ day per year, a day every $128$ years. The idealised Persian cycle uses $\frac8{33}\approx0.24242$, off by $0.00023$ day per year. This is a useful arithmetic model, not a complete description of the modern Iranian calendar, whose year begins at the astronomical equinox. The Gregorian $\frac{97}{400}=0.2425$ is not a convergent and is slightly worse than $\frac8{33}$, but it is convenient with a decimal century rule; the convergent $\frac{31}{128}=0.2421875$ would be far better still (off by $0.0000025$ day per year, a day in $400\,000$ years; long before that, the slow change of the year's length itself, a few tenths of a second per century, dominates).
\end{solution}

\begin{solution}{ex:n21}
\idea Run the algorithm on $\sqrt{n^2+1}$: the first complete quotient after the integer part is $\frac1{\sqrt{n^2+1}-n}=\sqrt{n^2+1}+n$ (rationalise), whose integer part is $2n$ and whose fractional part is the same as before. So the expansion repeats immediately. For the Pell solution, note $N(n+\sqrt{n^2+1})=n^2-(n^2+1)=-1$ and square.
\writeup Since $n<\sqrt{n^2+1}<n+1$, $a_0=n$ and $x_1=\frac1{\sqrt{n^2+1}-n}=\sqrt{n^2+1}+n$ (as $(\sqrt{n^2+1}-n)(\sqrt{n^2+1}+n)=1$). Its integer part is $2n$, and $x_1-2n=\sqrt{n^2+1}-n$, so $x_2=x_1$: $\sqrt{n^2+1}=[n;\overline{2n}]$. The number $u=n+\sqrt{n^2+1}$ has norm $n^2-(n^2+1)=-1$, so $u^2=(2n^2+1)+2n\sqrt{n^2+1}$ has norm $1$: $(2n^2+1)^2-(n^2+1)(2n)^2=1$.
\end{solution}

\begin{solution}{ex:sqtri}
\idea Complete the square: $\frac{m(m+1)}2=k^2$ becomes $(2m+1)^2-8k^2=1$, i.e.\ $x^2-2y^2=1$ with $x=2m+1$ odd and $y=2k$ even. Theorem~\ref{thm:pell} lists all solutions as powers of $3+2\sqrt2$; then check the parities.
\writeup Multiply by $8$ and add $1$: $\frac{m(m+1)}2=k^2$ iff $(2m+1)^2-2(2k)^2=1$. By Theorem~\ref{thm:pell} (fundamental solution $3+2\sqrt2$, see the example), the solutions with $x,y\ge1$ are $x_j+y_j\sqrt2=(3+2\sqrt2)^j$, i.e.\ $x_{j+1}=3x_j+4y_j$, $y_{j+1}=2x_j+3y_j$. By induction $x_j$ is odd and $y_j$ is even, so every positive solution gives $m=\frac{x_j-1}2$, $k=\frac{y_j}2$. The trivial solution $(x,y)=(1,0)$ adds the square triangular number $T_0=0$. The first ones are $(x,y)=(3,2),(17,12),(99,70),(577,408)$, giving $k^2=1,36,1225,41616$, i.e.\ $T_1=1$, $T_8=36$, $T_{49}=1225$, $T_{288}=41616$, and infinitely many more.
\end{solution}

\begin{solution}{ex:golden}
\idea The error formula of Theorem~\ref{thm:cfapprox} is exact: $q_n^2\left|x-\frac{p_n}{q_n}\right|=\frac1{x_{n+1}+q_{n-1}/q_n}$. For $\varphi$ every complete quotient is $\varphi$ and $\frac{q_{n-1}}{q_n}=\frac{F_n}{F_{n+1}}\to\frac1\varphi$; and $\varphi+\frac1\varphi=\sqrt5$. For the second part, any very good approximation is a convergent by Legendre, and we just computed how good the convergents are.
\writeup For $\varphi$, $x_{n+1}=\varphi$ and $q_n=F_{n+1}$, so $q_n^2\left|\varphi-\frac{p_n}{q_n}\right|=\frac1{\varphi+F_n/F_{n+1}}\to\frac1{\varphi+\varphi^{-1}}=\frac1{\sqrt5}$ (by (b) of Exercise~\ref{ex:cfunique}, $\frac{F_{n+1}}{F_n}\to\varphi$, and $\varphi-\varphi^{-1}=1$, $\varphi+\varphi^{-1}=\sqrt{(\varphi-\varphi^{-1})^2+4}=\sqrt5$). Let $c>\sqrt5$. If $\left|\varphi-\frac pq\right|<\frac1{cq^2}<\frac1{2q^2}$ with $\frac pq$ in lowest terms, then $\frac pq$ is a convergent (Theorem~\ref{thm:legendre}), $\frac{p_n}{q_n}$, with $q_n^2\left|\varphi-\frac{p_n}{q_n}\right|<\frac1c<\frac1{\sqrt5}$; by the limit, this happens for finitely many $n$ only.
\end{solution}

\begin{solution}{ex:liouville}
\idea (a) A rational $\frac pq$ near $x$ makes $P(\frac pq)$ small, since $P(x)=0$ and $P$ is Lipschitz near $x$; but $P(\frac pq)$ is an integer divided by $q^d$, and non-zero because an irreducible $P$ of degree $\ge2$ has no rational root. So $|P(\frac pq)|\ge\frac1{q^d}$, which forces $\frac pq$ to stay at distance $\gtrsim\frac1{q^d}$. (b) $L$ is approximated by its partial sums with an error far smaller than any $\frac c{q^d}$.
\writeup (a) Let $M=\sup_{[x-1,x+1]}|P'|$. If $\left|x-\frac pq\right|>1$ the bound holds with $c=1$. Otherwise, $q^dP(\frac pq)$ is an integer, non-zero since $P$ has no rational root (a rational root would give the factor $X-\frac pq$ in $\Q[X]$, and $P$ is irreducible of degree $\ge2$). So $\left|P(\frac pq)\right|\ge\frac1{q^d}$. By the mean value theorem, $\left|P(\frac pq)\right|=\left|P(\frac pq)-P(x)\right|\le M\left|x-\frac pq\right|$. Hence $\left|x-\frac pq\right|\ge\frac1{Mq^d}$, and $c=\min(1,\frac1M)$ works.

(b) Let $q_n=10^{n!}$ and $p_n=q_n\sum_{k\le n}10^{-k!}\in\Z$. Then $0<L-\frac{p_n}{q_n}=\sum_{k>n}10^{-k!}<2\cdot10^{-(n+1)!}=\frac2{q_n^{n+1}}$. $L$ is irrational: if $L=\frac ab$, then $L-\frac{p_n}{q_n}$ is a positive multiple of $\frac1{bq_n}$, so $\frac1{bq_n}<\frac2{q_n^{n+1}}$, false for large $n$. If $L$ were a root of a non-zero integer polynomial, it would be a root of one of its irreducible factors in $\Q[X]$, which can be taken in $\Z[X]$ (clear denominators), of some degree $d\ge2$ as $L$ is irrational; then (a) would give $\frac c{q_n^d}\le\frac2{q_n^{n+1}}$, i.e.\ $c\le2q_n^{d-n-1}\to0$: contradiction.
\end{solution}

\section*{Sources and further reading}

\begin{itemize}
\item A.~Ya.~Khinchin, \emph{Continued Fractions}, University of Chicago Press, 1964 (Dover reprint, 1997). A short classic, ending with the metric (probabilistic) theory of continued fractions: a first meeting of number theory and probability.
\item G.~H.~Hardy and E.~M.~Wright, \emph{An Introduction to the Theory of Numbers}, 6th ed., Oxford University Press, 2008, Chapters~10--11. Continued fractions, Hurwitz's theorem and Liouville numbers.
\item H.~W.~Lenstra Jr., Solving the Pell equation, \emph{Notices of the American Mathematical Society} 49 (2002), 182--192. History and algorithms, from Archimedes' cattle problem to modern methods.
\item J.~M.~Borwein, A.~van der Poorten, J.~Shallit and W.~Zudilin, \emph{Neverending Fractions}, Cambridge University Press, 2014. A modern tour, including the expansion of $e$ and computational aspects.
\item O.~Karpenkov, \emph{Geometry of Continued Fractions}, 2nd ed., Springer, 2022. The geometric view started by Klein and Arnold: convergents as vertices of the ``sails'', the convex hulls of the lattice points on each side of the line $p=xq$.
\item Archimedes' cattle problem: I.~Vardi, \emph{Archimedes' cattle problem}, American Mathematical Monthly 105 (1998), 305--319, with the epigram and its history; the full solution was first printed by H.~C.~Williams, R.~A.~German and C.~R.~Zarnke, Math.\ Comp.\ 19 (1965), 671--674.
\end{itemize}
